\chapter*{Executive Summary}
\addcontentsline{toc}{chapter}{Executive Summary}
\sloppy% lets the deployment URL wrap

\textbf{The problem.} When several AI agents cooperate on a task, a mistake made early can pass silently through later agents. In our running example, a refund workflow's order-lookup tool returns a stale price; every later agent behaves correctly on that input, and the customer is told a wrong refund amount in a fluent message, with no error anywhere. Developers today read logs by hand, usually blame the step where the error surfaced rather than where it began, and confirm a fix with a single re-run that a non-deterministic model can pass by chance.

\textbf{What we are building.} Granite is an evidence-first investigation workbench with the tagline "find the failure, prove the cause, verify the fix." It runs a workflow, checks its outcome and each step against executable rules and an independent source of truth, opens a finding when a completed run is wrong, ranks where the error began by following the data flow backwards, and verifies the top hypothesis by replaying the workflow five times with that step corrected and twice without, deciding CONFIRMED, REFUTED or INCONCLUSIVE under a pre-registered rule. A person then accepts or rejects the finding, and every artefact is stored as tamper-evident evidence and exported as a report. It runs at zero cost on local or recorded models, so sensitive traces need not leave the team.

\textbf{What changed after the proposal.} The panel asked for a better interface, clearer features and inputs, a concrete verification method and a redefined workflow. We found that all five comments came from one cause: we had designed mechanisms, not a user workflow. We redefined the product around one workflow and a "verified finding," replaced the read-only proof-of-concept console with a working application, and re-baselined our novelty after finding close prior work (DoVer and several commercial tools). Our contribution is now executable-check detection, deterministic provenance attribution, and a replay protocol with controls that is robust to model non-determinism, all self-hosted.

\textbf{Current state.} Iteration 1 is implemented and deployed at \url{https://granite-demo.onrender.com}. Eight features work end to end, 381 automated tests pass, continuous integration runs on every push, and 21 design diagrams are generated from the code. A pre-registered feasibility run on ten tasks met all six thresholds, caught a genuine model error that nobody injected, and exposed one design limit (two faults in a chain) that we have logged.

\textbf{Risks and gaps.} The implementation was produced in a short, AI-assisted burst committed from one account; each member now has to demonstrate ownership of their module and contribute through reviewed pull requests. Personal-data redaction before storage is not yet applied, the evaluation covers only one fault type and a small sample, and no usability test with outside users has been run.

\textbf{Remaining work.} Iteration 2 (to FYP-1 Final) adds trace import, a second workflow, three fault types, a Python SDK with a LangGraph adapter, fix suggestions, redaction, PDF reports and a first evaluation against baselines. FYP-2 adds a benchmark of at least 300 runs, a CrewAI adapter and full user acceptance testing.
